\documentclass[11pt]{article}
\usepackage{vmargin}
\setpapersize{A4}
\setmargins{15mm}{10mm}{180mm}{257mm}%changer 252 par 262
             {5mm}{2mm}{0pt}{10mm}
\usepackage[french]{babel}
\usepackage[utf8]{inputenc}

\usepackage{lmodern} 
\usepackage[T1]{fontenc}
\usepackage{fouriernc}
\usepackage{eurosym} % pour le symbole euro
\usepackage{graphicx} 
\usepackage{array} 
\usepackage{tabularx}
\usepackage{variations}
\usepackage{amsmath} 
\usepackage{caption}
\usepackage{amssymb}
\usepackage{awesomebox}
\usepackage{epic} % pour le dessin avec l'environnment picture
\usepackage{rotating}
\usepackage[autolanguage]{numprint}
\usepackage[tikz]{bclogo} 

\usepackage{circuitikz} % circuit electriques

\usepackage{pst-labo}
\usepackage{chemist} %chimie
\usepackage[version=4]{mhchem}


\setlength{\unitlength}{1 mm}
\usepackage{xy} % pour la réalisation de dessin, graphique, diagramme
\xyoption{all}

\newcommand{\cadreblanc}[1]{%
\medskip\framebox[\linewidth][c]{%
\rule{0mm}{#1mm}}\par
\medskip}


\def\labelenumii{\theenumi.~\theenumii.}
\def\theenumii{\arabic{enumii}}


\presetkeys{bclogo}{logo=\bcquestion,couleur=gray!10,barre=snake,arrondi = 0.1}{}

\newcommand{\ctext}[1]
{\begin{bclogo}[couleur = white,logo= , arrondi = 0.1,barre=none]
{}
\vspace{-8mm}
~~

~~
\end{bclogo}}

  
\newcommand{\defbox}{\awesomebox[black]{2pt}{\fbox{
\includegraphics[scale=0.5]{Chapitre 1 - Energie et puissance électrique/images/appel exam.png}}}{black}}

%%%%%%%%%%%%%%%%%%%%%%%%%%%%%%%%%%%%%%%%%%%%%%%%%%%%%%%%%%
%
% Définition des sections et sous-sections
%
%%%%%%%%%%%%%%%%%%%%%%%%%%%%%%%%%%%%%%%%%%%%%%%%%%%%%%%%%%


\renewcommand{\thesection}{{\arabic{section}.}}

\addto\captionsfrench{\def\tablename{Document}}
\usepackage{setspace}
\usepackage{multicol}
\usepackage{multirow}
\usepackage{fancyhdr}
\usepackage{graphpap}

%%%%%%%%%%%%%%%%%%%%%%%%%%%%%%%%%%%%%%%%%%%%%%%%%%%%%%%%%%
%
% Définition de nouveaux types de colonne
%
%%%%%%%%%%%%%%%%%%%%%%%%%%%%%%%%%%%%%%%%%%%%%%%%%%%%%%%%%%

\newcolumntype{C}{>{\centering\arraybackslash}X}
\newcolumntype{L}{>{\raggedright\arraybackslash}X}
\newcolumntype{R}{>{\raggedleft\arraybackslash}X}


\renewcommand{\footrulewidth}{0.4pt}
\usepackage{lastpage} 

\newcommand{\appel}{\raisebox{-2ex}{
\includegraphics[width=1cm]{images/appel.jpg}}\quad}



%%%%%%%%%%%%%%%%%%%%%%%%%%%%%%%%%%%%%%%%%%%%%%%%%%%%%%%%%%
%                                                        %
%              Evaluation par compétences                %
%                                                        %
%%%%%%%%%%%%%%%%%%%%%%%%%%%%%%%%%%%%%%%%%%%%%%%%%%%%%%%%%%



\newcommand\SA{\marginpar{\scalebox{.8}{\framebox[1.6ex]{\rule{0ex}{1.6ex}\raisebox{.75ex}{\rule{1.6ex}{.5pt}}}\hspace{-.3pt}\framebox[3.2ex]{S}}}}
\newcommand\A{\marginpar{\scalebox{.8}{\framebox[1.6ex]{\rule{0ex}{1.6ex}\raisebox{.75ex}{\rule{1.6ex}{.5pt}}}\hspace{-.3pt}\framebox[3.2ex]{A}}}}
\newcommand\R{\marginpar{\scalebox{.8}{\framebox[1.6ex]{\rule{0ex}{1.6ex}\raisebox{.75ex}{\rule{1.6ex}{.5pt}}}\hspace{-.3pt}\framebox[3.2ex]{R}}}}
\newcommand\V{\marginpar{\scalebox{.8}{\framebox[1.6ex]{\rule{0ex}{1.6ex}\raisebox{.75ex}{\rule{1.6ex}{.5pt}}}\hspace{-.3pt}\framebox[3.2ex]{V}}}}
\newcommand\C{\marginpar{\scalebox{.8}{\framebox[1.6ex]{\rule{0ex}{1.6ex}\raisebox{.75ex}{\rule{1.6ex}{.5pt}}}\hspace{-.3pt}\framebox[3.2ex]{C}}}}


%%%%%%%%%%%%%%%%%%%%%%%%%%%%%%%%%%%%%%%%%%%%%%%%%%%%%%%%%%
%                                                        %
% Définitions de l'année, la spécialité, du diplome ...  %
%                                                        %
%%%%%%%%%%%%%%%%%%%%%%%%%%%%%%%%%%%%%%%%%%%%%%%%%%%%%%%%%%


\newcommand{\specialite}{Cybersécurité, Informatique
et Réseaux, Électronique (CIEL)}
\newcommand{\annee}{2024-25}
\newcommand{\diplome}{Baccalauréat Professionnel}
\newcommand{\etablissement}{LP Denis Diderot}
\newcommand{\classe}{1CIEL1}
\newcommand{\datedeval}{27 Mai 2025}
\newcommand{\matiere}{Sciences Physiques et chimiques }

\newcommand{\sequence}{1/2}
\newcommand{\theme}{
Caractériser quantitativement une solution aqueuse}
%%%%%%%%%%%%%%%%%%%%%%%%%%%%%%%%%%%%%%%%%%%%%%%%%%%%%%%%%%
%                                                        %
% Définitions des capacités, connaissances et attitudes. %
%                                                        %
%%%%%%%%%%%%%%%%%%%%%%%%%%%%%%%%%%%%%%%%%%%%%%%%%%%%%%%%%%





\newcommand{\capacites}{
  \begin{minipage}[c]{12.5cm}
    \begin{itemize}
      \item \rule{0ex}{2.5ex} Réaliser une solution de concentration en quantité de matière donnée par dilution ou dissolution. 
      \item  Calculer une masse molaire moléculaire à partir des masses molaires atomiques et de la formule chimique de la molécule.  
      \item Calculer la concentration en masse d’un soluté à partir de sa concentration en quantité de matière et de sa masse molaire moléculaire.  \rule[-1.5ex]{0ex}{3ex} 
      \end{itemize}
    \end{minipage}
  }




\newcommand{\connaissances}{
  \begin{minipage}[c]{12.5cm}
    \begin{itemize}
    \item \rule{0ex}{2.5ex}Connaître les définitions d'une solution, d'un solvant, d'un soluté.
    \item Connaître les relations entre la masse molaire $n$, la masse d’un échantillon $m$ et la quantité de matière $n$ 
    \item Connaître la relation entre la concentration en quantité de matière de soluté $c$, la quantité de matière $n$ et le volume de la solution $V$ 
    \item Connaître la définition de la concentration en masse d’un soluté dans une solution.  \rule[-1.5ex]{0ex}{3ex}
    \end{itemize}
    \end{minipage}
  }

\newcommand{\attitudes}{
  \begin{minipage}[c]{12.5cm}
    \begin{itemize}
    \item \rule{0ex}{2.5ex}Le sens de l'observation.
    \item L'ouverture à la communication, au dialogue et au débat argumenté.
    \item La rigueur et la précision
    \item Le respect des règles élémentaires de sécurité. 
    \item \rule[-1.5ex]{0ex}{3ex}Le gout de chercher et de raisonner.
    \end{itemize}
  \end{minipage}
}




%%%%%%%%%%%%%%%%%%%%%%%%%%%%%%%%%%%%%%%%%%%%%%%%%%%%%%%%%%
%                                                        %
%                   Début du document                    %
%                                                        %
%%%%%%%%%%%%%%%%%%%%%%%%%%%%%%%%%%%%%%%%%%%%%%%%%%%%%%%%%%





%%%%%%%%%%%%%%%%%%%%%%%%%%%%%%%%%%%%%%%%%%%%%%%%%%%%%%%%%%
%                                                        %
%                     Page de garde                      %
%                                                        %
%%%%%%%%%%%%%%%%%%%%%%%%%%%%%%%%%%%%%%%%%%%%%%%%%%%%%%%%%%



\begin{document}
\pagestyle{fancy}
\lfoot{Alexis Ruiz}
\cfoot{\thepage/3}
\rfoot{\datedeval}

\thispagestyle{empty}
\Large 
\noindent\begin{tabularx}{1\linewidth}[c]{|c|c|c|C|C|}
  \hline
  \multirow{2}{*}{
  \begin{minipage}[c]{3cm}
    
\includegraphics[width=30mm]{images/aca-besancon.png}
  \end{minipage}}   &
 \multicolumn{4}{c|}{
 \begin{minipage}[c]{12cm}
      \begin{center}
         \rule[-1.5ex]{0ex}{5ex}BACCALAURÉAT PROFESSIONNEL\\
         \rule[-2ex]{0ex}{4ex} \specialite 
      \end{center}    \end{minipage}}\\
  \cline{2-5}
  & \multicolumn{3}{c|}{\rule[-2ex]{0ex}{5.5ex}Épreuve E1- Sciences Physiques et chimiques } & Coef. 2\\
  \hline
  \multirow{2}{*}{
  \begin{minipage}[c]{3cm}
    \begin{center}
      Contrôle en \\ Cours de\\ Formation
    \end{center}      
    \end{minipage}}
    &
  \multirow{2}{*}{
    \raisebox{-2ex}{\begin{minipage}[c]{6.3cm}
      \begin{center}
        Situation d'évaluation de\\ \matiere
     \end{center}  
    \end{minipage}}}
    &
  \multirow{2}{*}{
    \begin{minipage}[c]{2.2cm}
      \begin{center}
        Année\\ scolaire\\ \annee
  \end{center}  
    \end{minipage}}
    &
    \rule[-1ex]{0ex}{4ex}Séquence &
    Durée\\
    \cline{4-5}
    &&&
    \rule[-1.5ex]{0ex}{5ex}\sequence
    & 
    45~min
\\
\hline
\end{tabularx}

\bigskip

\noindent\fbox{
  \begin{minipage}[c]{.97\linewidth}
    \centerline{\rule[-.8ex]{0ex}{3.5ex}FICHE D'INFORMATION DU CANDIDAT}
  \end{minipage}
}

\bigskip


\normalsize

\noindent\begin{tabularx}{1\linewidth}{|p{11cm}C|}
  \hline
  \rule{0ex}{3ex}Établissement~:~\textbf{{\etablissement}} & Classe~:~\textbf{\classe}\\
 \rule[-1ex]{0ex}{4ex} NOM et Prénom du CANDIDAT : \ \dotfill & Date de l'évaluation~:~\textbf{\datedeval}\\
 \hline
\end{tabularx}
\bigskip

\noindent{\Large Thème~: \theme}\\

\noindent{\Large Module(s) évalué(s)~: }

\begin{center}
  Capacités, connaissances et attitudes du référentiel évaluées\\
  \begin{tabularx}{16cm}{|C|l|}
    \hline
    \textbf{Capacités} &
    \capacites\\
    \hline 
    \textbf{Connaissances} &
    \connaissances    \\
    \hline
    \textbf{Attitudes}  &
    \attitudes\\
    \hline
  \end{tabularx}
\end{center}

\noindent\fbox{\begin{minipage}[c]{.97\linewidth}
    \begin{description}
  \item[$\vartriangleright$] La clarté des raisonnements et la qualité de la rédaction
    interviendront dans l'appréciation des copies.
  \item[$\vartriangleright$] L'emploi des calculatrices est autorisé, dans les conditions
    prévues par la réglementation en vigueur.
  \end{description}
\end{minipage}}

\bigskip


L'évaluation est notée sur 10 : 7 points pour la mobilisation des
connaissances et des compétences pour résoudre des problèmes et 3
points pour l'activité d'expérimentation.

\begin{center}
  \begin{tabularx}{1\linewidth}{|L|c|}
  \hline
  \textbf{\rule{0ex}{3ex}Le candidat atteste avoir été informé de la date et des objectifs
    de l'évaluation le \makebox[4cm]{\rule{0ex}{3ex}\dotfill}} & \textbf{\rule[-6ex]{0ex}{3ex}\underline{Émargement}}\\
   \hline
\end{tabularx}
\end{center}




%%%%%%%%%%%%%%%%%%%%%%%%%%%%%%%%%%%%%%%%%%%%%%%%%%%%%%%%%%
%                                                        %
%                   fin page de garde                    %
%                                                        %
%%%%%%%%%%%%%%%%%%%%%%%%%%%%%%%%%%%%%%%%%%%%%%%%%%%%%%%%%%


\newpage

\lhead{\textbf{\matiere}}
\rhead{Année scolaire \annee}
\setcounter{page}{1}
\large

\noindent\begin{tabularx}{1\linewidth}[c]{|c|c|C|C|}
  \hline
  \multirow{2}{*}{\Large
    \begin{minipage}[c]{4cm}
      \begin{center}
        Contrôle en Cours\\ de Formation
    \end{center}      
  \end{minipage}}
    &
    \multirow{2}{*}{
      \begin{minipage}[c]{7.2cm}
        \begin{center}
          Situation de Sciences Physiques et chimiques
        \end{center}  
    \end{minipage}}
  &
  \rule[-1ex]{0ex}{4ex}Séquence &
  Durée\\
  \cline{3-4}
  &&
    \rule[-1.5ex]{0ex}{5ex}\sequence
    & 
    45~min\\
\hline
\end{tabularx}
\medskip


\noindent\fbox{
  \Large\begin{minipage}[c]{.97\linewidth}
    \centerline{\rule[-.8ex]{0ex}{3.5ex}SUJET DESTINÉ AU CANDIDAT}
  \end{minipage}
}

\medskip

\noindent\begin{tabularx}{1\linewidth}{|p{10cm}L|}
  \hline
  \rule{0ex}{3ex}Établissement:~\textbf{{\etablissement}} & Classe~:~\textbf{\classe}\\
 \rule[-1ex]{0ex}{4ex} NOM et Prénom : \ \dotfill & Date de l'évaluation : \textbf{\datedeval}\\
 \hline
\end{tabularx}

\medskip


\noindent\begin{tabularx}{1\linewidth}{|c|L|}
  \hline
\multirow{2}{1.3cm}{ \appel} & \textbf{\rule[-.8ex]{0ex}{3.5ex}L'examinateur intervient à la demande du
candidat ou quand il le juge utile.}\\
\cline{2-2} 
&\textbf{\rule[-.8ex]{0ex}{3.5ex}Dans la suite du document, ce symbole  \raisebox{-1.2
ex}{
\includegraphics[scale=0.35]{Chapitre 1 - Energie et puissance électrique/images/appel exam.png}}~ signifie <<Appeler
  l'examinateur>>.}\\
 \hline
\end{tabularx}

\bigskip

\noindent
\normalsize
\setstretch{1.2}

\parindent 0cm


\noindent
  \setlength{\arraycolsep}{2pt}
  \begin{bclogo}[logo=\bcinfo]
{Situation~: Lors d'un effort prolongé, un marathon par exemple, le sportif perd de l'eau et des sels minéraux en transpirant et puise dans ses réserves énergétiques pour alimenter ses muscles en énergie. De nombreuses boissons sont proposées dans le commerce pour compenser ces pertes (boisson isotonique).   }

\end{bclogo}

\begin{bclogo}{~Problématique :}

    \textbf{Après un marathon, Kevin, un sportif présente une fatigue extrême. Arrivé chez lui, il constate qu'il n'a plus de boisson isotonique. Il décide de fabriquer lui-même 100 mL d'une solution lui permettant une récupération rapide à partir de chlorure de
sodium (sel).} \\[2mm]
Il demande à une IA, la concentration molaire C en mol/L en chlorure de sodium (NaCl) nécessaire pour fabriquer un volume de  100 mL cette boisson isotonique. 

\begin{center}
    
\fbox{


\begin{minipage}{14cm}
\vspace{3mm}
 Chlorure de Sodium \hfill  Formule chimique : NaCl \hfill  C = 0,085 mol/L
 \vspace{3mm}
\end{minipage}
}
\end{center}

\vspace{2mm}



\end{bclogo}


 \begin{enumerate}
\item \textbf{Compréhension et analyse de la situation / formulation d'une hypothèse.}\\

\begin{enumerate}


\item On donne ci-dessous un extrait de classification périodique ; montrer que la masse molaire du chlorure de sodium \textbf{M($\mathrm{NaCl}$)} est égale à 58,5 g/mol. 

\begin{center}
\begin{tikzpicture}
\draw (0,0)--(0,3)--(2,3)--(2,0)--cycle;
\node at (1,1.75) {{\LARGE Na}};
\node at (1,0.3) {23};
\node at (1,1) {Sodium};
\node at (0.3,2.6){{\Large 11}};
\end{tikzpicture}
\hspace{1cm}
\begin{tikzpicture}
\draw (0,0)--(0,3)--(2,3)--(2,0)--cycle;
\node at (1,1.75) {{\LARGE Cl}};
\node at (1,0.3) {35,5};
\node at (1,1) {Chlore};
\node at (0.3,2.6){{\Large 17}};
\end{tikzpicture}
\end{center}


\begin{minipage}[c]{.91\linewidth}

    \cadreblanc{15}
    
  \end{minipage}\SA\R
\vspace{3mm}

\newpage

\item Retrouver dans la problématique la concentration en chlorure de sodium C, en mol/L, préconisée par l'IA. \SA
 \C
 
\vspace{8mm}
\begin{minipage}[c]{.91\linewidth}

    .\dotfill\\
    
  \end{minipage}

  \vfill 

\item Calculer, en mol, \textbf{n} la quantité de matière en chlorure de sodium contenue dans un volume de 100 mL de solution . On utilisera la formule \fbox{$\mathrm{n = C \times V }$}

\begin{minipage}[c]{.91\linewidth}

    \cadreblanc{10}
    
  \end{minipage}\R
\vfill 

\item Calculer, en g, \textbf{m} la masse de NaCl nécessaire à la préparation de la solution.  \\ On utilisera la formule \fbox{$\mathrm{n=\dfrac{m}{M}}$}. \\ 
\texbf{ Arrondir le résultat à 0,1 g. }

\begin{minipage}[c]{.91\linewidth}

    \cadreblanc{10}
    
  \end{minipage}\R
\vfill 

\item Kevin possède une balance précise à 0,1 g près, pourra-t-il effectuer cette pesée avec précision ? Justifier votre réponse. 

\begin{minipage}[c]{.91\linewidth}

    \cadreblanc{10}
    
  \end{minipage}\V \C
\end{enumerate}

\vfill 

\textbf{\item Expérimentation }
\begin{enumerate}

\item Remettez dans l'ordre les différentes 4 étapes d'une dissolution.    \\[5mm]
\begin{minipage}[c]{.91\linewidth}

\includegraphics[scale=0.73]{Chapitre 4 - Solutions aqueuses/images/disolutionCCF.png}
    
  \end{minipage}\A



\defbox{\fbox{
\begin{minipage}{12cm}{

\textbf{Appel $\mathrm{n^\circ}$1 :} Appeler l'examinateur  pour qu'il vérifie vos calcul et l'ordre des étapes  puis réaliser la dissolution \fbox{avec une masse \textbf{m = 5g} }}

\end{minipage}
}}  \R

\vfill 

\newpage

\item Calculer, en mol/L, la concentration $\mathrm{C_0}$ de la solution. \R


\begin{minipage}[c]{.91\linewidth}

    \cadreblanc{10}
    
  \end{minipage}
\vspace{3mm}

\item Cette solution est-elle en accord avec la valeur préconisée par l'IA ? \\[0.5cm]

\begin{minipage}[c]{.91\linewidth}

    .\dotfill\\
    
  \end{minipage}\V
\vspace{3mm}




\vfill 




    \item On souhaite dilué cette solution 10 fois, choisir dans la liste ci-dessous le matériel nécessaire :  \\[5mm]
    \SA \A 
\vspace{-10mm}

    \begin{multicols}{3}

    $\square$ Bécher \\
    $\square$ Entonnoir \\
    $\square$ Eau distillée 
    \columnbreak

    
    $\square$ Pipette jaugée 5 mL \\
    $\square$ Pipette jaugée 10 mL \\
    $\square$ Pipette jaugée 20 mL 

        \columnbreak


    $\square$ Fiole jaugée 50 mL \\
    $\square$ Fiole jaugée 100 mL  \\
    $\square$ Fiole jaugée 250 mL 

    \end{multicols}

    \vfill

    \defbox{\fbox{
\begin{minipage}{12cm}{

\textbf{Appel $\mathrm{n^\circ}$2 :} Appeler l'examinateur  pour qu'il vérifie votre choix de matériel. \textbf{Puis
réaliser la dilution devant lui} en suivant le protocole ci-dessous. }

\end{minipage}
}} \R

\end{enumerate}

\vfill 

\includegraphics[scale=0.88]{Chapitre 4 - Solutions aqueuses/images/dilution3.png}

\vfill 

\textbf{\item Exploitation et conclusion }
\begin{enumerate}

\item Calculer, en mol/L, la concentration $\mathrm  C$ de la solution diluée. 

  \vspace{3mm}

\begin{minipage}[c]{.91\linewidth}

\cadreblanc{20}
    
  \end{minipage}\R
  



    
\vfill 


  \item Cette valeur est-elle en accord avec la concentration de la solution suggérée par l'IA ? \\ 



\begin{minipage}[c]{.91\linewidth}

\vspace{8mm}
    .\dotfill\\
    
  \end{minipage}\V \C

  




\end{enumerate}

\end{enumerate}

\vfill 

\newpage

%%%%%%%%%%%%%%%%%%%%%%%%%%%%%%%%%%%%%%%%%%%%%%
%%%%%%%%%%%%%%%%%%%%%%%%%%%%%%%%%%%%%%%%%%%%%%
%%                                          %%
%%          Fiche d'évaluation              %%
%%                                          %%
%%%%%%%%%%%%%%%%%%%%%%%%%%%%%%%%%%%%%%%%%%%%%%
%%%%%%%%%%%%%%%%%%%%%%%%%%%%%%%%%%%%%%%%%%%%%%

\small

\thispagestyle{empty}
\setstretch{1}

\noindent\begin{tabularx}{1\linewidth}[c]{|L|l|l|}
  \hline
  \multicolumn{3}{|c|}{\large
    \begin{minipage}[c]{.7\linewidth}
      \begin{center}
        \bfseries \rule{0ex}{3ex}GRILLE NATIONALE D'ÉVALUATION

        EN MATHÉMATIQUES ET EN

        SCIENCES PHYSIQUES ET CHIMIQUES\rule[-1ex]{0ex}{3ex}
      \end{center}
    \end{minipage}
}\\
  \hline
  \rule[-1ex]{0ex}{4ex} Nom : & \rule{.3cm}{0cm}Diplôme
  préparé : \diplome\rule{.3cm}{0cm} &  Séquence
  d'évaluation\footnote{Chaque séquence propose la résolution de 
    problèmes issus du domaine professionnel ou de la vie courante. En
    mathématiques, elle comporte un ou deux exercices ; la résolution
    de l'un d'eux nécessite la mise en \oe{}uvre de capacités
    expérimentales.} \no\sequence \\
 \hline\end{tabularx}\\[.3cm]


\textbf{1.  Liste des capacités, connaissances et attitudes évaluées} 
\begin{center}
  \begin{tabularx}{1\linewidth}{|C|l|}
    \hline
    \textbf{Capacités} &
    \capacites\\
    \hline 
    \textbf{Connaissances} &
    \connaissances    \\
    \hline
    \textbf{Attitudes}  &
    \attitudes 
    \\
    \hline
  \end{tabularx}
\end{center}


\medskip

\textbf{2. Évaluation\footnote{
    {\bfseries En sciences physiques et chimiques :}  L'évaluation porte
    nécessairement sur des capacités expérimentales. 3 points sur
    10 sont consacrés aux questions faisant appel à la compétence
    \og{}Communiquer\fg{}.}}
\begin{center}
\begin{tabularx}{.98\linewidth}[c]{|c|C|c|p{0.42cm}|p{0.42cm}|p{0.4cm}|c|}
  \hline
  \textbf{Compétences} & 
  \textbf{Capacités}&\textbf{Questions} &\multicolumn{4}{c|}{
  \begin{minipage}[c]{2.8cm}
    \begin{center}
      \textbf{\rule{0ex}{2.2ex}Appréciation\\ du niveau\\
        d'acquisition\rule[-1.2ex]{0ex}{2ex}} 
    \end{center}    \end{minipage}\footnote{Le professeur peut utiliser toute forme d'annotation lui permettant d'évaluer l'élève (le candidat) par compétences.}}
    \\
    \hline
    S'approprier&
    Rechercher, extraire et organiser l'information.
    &\begin{minipage}[c]{1.2cm}
      \begin{center}
        \rule{0ex}{2.2ex}
        1.1 ; 1.2 ; \\ 2.4
\rule[-1.2ex]{0ex}{2ex}
      \end{center}\end{minipage}
    &&& & \\
    \hline    
    Analyser& 
    \begin{minipage}{1.0\linewidth}
  \begin{itemize}
  \item Émettre une conjecture, une hypothèse.
  \item Proposer une méthode de résolution, un protocole expérimental.
  \end{itemize}
      \end{minipage}
    &\begin{minipage}[c]{1.2cm}
      \begin{center}
        \rule{0ex}{2.2ex} 2.1 ; 2.4
        \rule[-1.2ex]{0ex}{2ex}
      \end{center}
    \end{minipage}
    & & & & \\
    \hline
    Réaliser& 
    \begin{minipage}{1.0\linewidth}
      \begin{itemize}
      \item \rule{0ex}{2.2ex}Choisir une méthode de résolution, un
        protocole expérimental.
      \item Exécuter une méthode de
        résolution, expérimenter, simuler.\rule[-1.2ex]{0ex}{2ex}
      \end{itemize}
    \end{minipage}&
    \begin{minipage}[c]{1.2cm}
      \begin{center}
    \rule{0ex}{2.2ex}
1.1 ; 1.3 ; 1.4; \textbf{A1}; 2.2 ; 3.1 ;\\ \textbf{A2} \rule[-1.2ex]{0ex}{2ex}
  \end{center}
    \end{minipage}
    
    & & & & \\
\hline
Valider &\begin{minipage}{1.0\linewidth}
  \begin{itemize}
  \item \rule{0ex}{2.2ex}Contrôler la vraisemblance d'une conjecture, d'une hypothèse.
  \item Critiquer un résultat, argumenter.\rule[-1.2ex]{0ex}{2ex}
  \end{itemize}
      \end{minipage}&
    \begin{minipage}[c]{1.2cm}
      \begin{center}
    \rule{0ex}{2.2ex}
1.5 ; 2.3 ; \\ 3.2
\rule[-1.2ex]{0ex}{2ex}
      \end{center}
    \end{minipage}
    & & & & \\
    \hline
   Communiquer &\begin{minipage}{1.0\linewidth}
   \begin{itemize}
       \item \rule{0ex}{2.2ex}Rendre compte d'une démarche, d'un résultat,
    à l'oral ou à l'écrit.\rule[-1.2ex]{0ex}{2ex}
   \end{itemize}
    \end{minipage}
    &
    \begin{minipage}[c]{1.2cm}
      \begin{center}
    
1.2 ; 1.5 ;\\ 3.2
      \end{center}
      \end{minipage}
      & & & &  \\
      \hline
      \multicolumn{3}{c|}{}& \multicolumn{4}{r|}{\rule[-1ex]{0ex}{4ex}\rule{3ex}{0ex}{\Large /~10}}\\    
    \cline{4-7}
  \end{tabularx}
\end{center}

\end{document}
